% --- Frontmatter ---
\thispagestyle{empty}
\begin{titlepage}
  \begin{center}
    \vspace*{1.5cm}
    {\Large\textbf{Trabajo de investigacion}}\\[0.6cm]
    {\Large Aplicado a la deteccion temprana de fallas}\\
    {\Large en subestaciones electricas de distribucion}\\[1.5cm]

    \rule{0.85\textwidth}{0.5pt}\\[0.6cm]
    {\Huge\bfseries Gemelo Digital ligero\\ de Subestaci\\'on}\\[0.3cm]
    {\Large Deteccion temprana no supervisada con SCADA/PMU,\\
        alertas calibradas y aprendizaje federado entre subestaciones}\\[0.6cm]
    \rule{0.85\textwidth}{0.5pt}\\[1.5cm]

    {\Large Miguel Ortiz Coilla}\\[1cm]
    {\large Octubre 2026}\\[2cm]

    \begin{tcolorbox}[colback=blue!3,colframe=blue!40!black,width=0.78\textwidth]
      \centering\itshape\small
      Las subestaciones de distribucion generan streams continuos de SCADA y PMU,
      pero las fallas son raras y los datos no siempre pueden salir del recinto.
      Este trabajo propone un pipeline de deteccion no supervisada basado en un
      gemelo digital ligero, modelos autoencoder LSTM y \emph{conformal prediction},
      entrenado en forma federada entre subestaciones para preservar la privacidad
      de los datos y cumplir con la normativa de ciberseguridad del sector
      electico chileno.
    \end{tcolorbox}
  \end{center}
\end{titlepage}